% page 403 of the facsimile (image p-402 of the scan)
% block 165.1 x 237.7 mm ; left margin 36.9 mm
\pgc{15.240mm}{0.000mm}{
\l{\cel{55}395}
\l{}
\l{nitro. (kemio). Nitrato naturala di kalio, NO3K.}
\l{}
\l{nivar. (netrans.) (aludante la aquo qua konjelesis en la supra}
\l{\cel{3}\sou{regioni di atmosfero}). Falar adsule, forme di floki blanka.IS.}
\l{}
\l{\sou{niuzo.}\cel{3}(\sou{tekn.)}\cel{2}(\sou{aludante jurnali, inform-agenterii}): Sorto}
\l{\cel{3}specala di informi : redakturo qua naracas fakti ne-expektita,}
\l{\cel{3}di qui la exakteso esas verifikita rigoroze (quale agus lo}
\l{\cel{3}historiisto), e qua eliminas la detali qui ne interesus la}
\l{\cel{3}lektonti. - E.}
\l{}
\l{\sou{nivelo}. Grado di alteso, kompare a plano horizontala, dilineo}
\l{\cel{3}o di plano qua esas paralela ad olu. - DFIRS.}
\l{}
\l{\sou{niveolo}. (\sou{bot.)}\cel{3}Genero de amarilidei, planto bulboza, di qua}
\l{\cel{3}la folii esas lineala e la flori blanka. L.\cel{1}\sou{galanthus nivalis}.}
\l{}
\l{\sou{nixo}. (\sou{mitol.,Germana e Skandinava}). Feo qua habitas profunde}
\l{\cel{3}en la aquo amasi. - DE.}
\l{}
\l{\sou{no.}\cel{3}Interjeciono uzata por enuncar ke kozo ne esas tala. - EFIS}
\l{\cel{3}(Se la adverso insistas pri la \textquotedbl{}eso\textquotedbl{},lu uzas la interjecioni}
\l{\cel{4}\sou{\textquotedbl{}si\textquotedbl{}}). -EFIS.}
\l{}
\l{\sou{nobelo}. Persono qua apartenas a klaso quan onu egardas kom supe-}
\l{\cel{3}riora, e qua honorizesis per privileji e distingivi heredita.-}
\l{\cel{3}DEFIS.}
\l{}
\l{\sou{nobla}. I.Elevita, ek la ceteri, per digneso o merito. - II.}
\l{\cel{3}Elevita super to quo esas ordinara. DEFIS.}
\l{}
\l{\sou{nocar}.\cel{2}(\sou{trans., pri, pro, per}). Produktar detrimento (ad ulu)}
\l{\cel{3}o domajo (ad ulo). - EFI.}
\l{}
\l{\sou{nocho}. I. Entaliuro quan facas panvendisto, karnovendisto, sur}
\l{\cel{3}bastono por indikar la quanto de pano, de karno, vendita}
\l{\cel{3}kredite. - II. Entaliuro sur la fusto di arbalesto o sur la}
\l{\cel{3}grosa extremajo di flecho, ube esas la kordo tensita di la}
\l{\cel{3}arbalesto,di la arm-arko. - E.}
\l{}
\l{\sou{nociono.}\cel{1}Savajo elementala. - DEFIS.}
\l{}
\l{\sou{nodo}. Krucumuro tre klemita qua formacas halto en la kontinueso}
\l{\cel{3}di filo, di kordo, e c., o qua unionas strete du fili, du}
\l{\cel{3}kordi, e c.\cel{2}- FIS.}
\l{}
\l{noktar.\cel{2}(\sou{netrans.)}\cel{2}Dicesas pri la parto di la dio dum qua la}
\l{\cel{3}suno-radii cesis lumizar la loko ube onu esas. - DEFIRS.}
\l{}
\l{noktiluko. (\sou{zool.)}\cel{3}Genero de protozoo flogilatra, ek la klaso}
\l{\cel{3}\textquotedbl{}cistoflagelei\textquotedbl{} - di qua la \textquotedbl{}flogilumo\textquotedbl{}, tre povoza, atingas}
\l{\cel{3}un milimetro diametre kande lu esas adulta, fosforecanta.}
\l{\cel{3}L. noctiluca.\cel{2}EFISL.}
}
